\documentclass[conference]{IEEEtran}
\IEEEoverridecommandlockouts

\usepackage{amsmath,amssymb}
\usepackage{booktabs}
\usepackage{cite}
\usepackage{graphicx}
\usepackage{url}
\usepackage{xcolor}

\newcommand{\thz}{THZ1}
\newcommand{\thp}{THP1}

\title{Joint Dwell-Time and Charging-Power Control of Multi-Line Catenary-Free Tramways for Shared Demand Management}

\author{
\IEEEauthorblockN{Jun-Jie Zhao}
\IEEEauthorblockA{School of Electrical and Computer Engineering\\
The University of Sydney\\
Sydney, Australia}
}

\begin{document}
\maketitle

\begin{abstract}
Urban rail operators are seeking cleaner and more flexible ways to power and run their services. Catenary-free tramways avoid continuous overhead wires by using onboard energy storage and rapid charging at stations. This design supports sustainable transport and reduces the visual impact of rail infrastructure. However, most existing studies consider only one line and do not address how to coordinate several catenary-free tram lines whose vehicles arrive at different times. When arrivals and charging periods overlap, the operator's total power demand can rise sharply. Stop duration also affects charging time, passenger boarding and alighting, and later vehicle movements. Multi-line operation must therefore balance total power demand, sufficient onboard energy, and service efficiency in real time. This paper presents the first study of joint stop-duration and charging-power control for multi-line catenary-free tramways and proposes an event-driven deep reinforcement learning method. A discrete-event simulation models two lines with different onboard storage technologies. At each arrival, the online decision-making agent selects the stop duration and charging power under a shared power-demand budget. A case study based on Guangzhou Haizhu Tram Line~1 and Huangpu Tram Line~1 shows that the method reduces simultaneous charging peaks and budget exceedance while keeping enough energy onboard and maintaining passenger service under several running-time disturbances and budget settings. It performs better overall than fixed operation and the genetic algorithm, particle swarm optimization, differential evolution, and model predictive control baselines in managing total demand and making real-time decisions. The method is therefore a practical option for coordinating multiple catenary-free tram lines in changing operating conditions.
\end{abstract}

\begin{IEEEkeywords}
Catenary-free tramway, electric public transport, demand management, reinforcement learning, discrete-event simulation, proximal policy optimization.
\end{IEEEkeywords}

\section{Introduction}
Catenary-free light rail and tram systems provide electric public transport without continuous overhead wiring by using onboard energy storage and station-based charging. Removing much of the catenary, support poles, and related equipment reduces visual intrusion and improves urban aesthetics, a particularly valuable benefit in dense city centers and historic districts \cite{guerrieri2019catenary}. Catenary-free light rail and tram recharge their onboard storage during short station stops. Supercapacitors are well suited to this duty because they can accept high charging power, but this capability also transfers a large, short-duration load to the station power supply. If several vehicles charge simultaneously, their demands add and may exceed the capacity planned for feeders and transformers or the operator's contracted demand, creating equipment-loading and voltage concerns or costly demand peaks. In this paper, the \emph{total charging-power budget} is the operator's planning threshold for simultaneous aggregate grid-side charging power, measured in kW. Each arrival must receive enough charging time and power for subsequent sections, while passengers need time to board and alight. More charging power reduces the risk of an energy shortfall but consumes more of the shared limit, whereas a longer stop can disturb vehicle headways---the time gaps between consecutive vehicles---and increase passenger waiting.

Research relevant to this problem spans catenary-free railway energy management, railway timetabling and rescheduling, opportunity charging, in which vehicles recharge during scheduled service stops, and feedback and learning control. Catenary-free tramways rely on onboard storage and repeated high-power charging, so their operating limits depend on both storage technology and service conditions \cite{guerrieri2019catenary}. Previous work has optimized speed profiles for battery-equipped catenary-free trains \cite{ghaviha2019speed}, but it does not coordinate stop duration and charging power across several services that share one operator budget.

% Preserve the prose citation order although the double-column comparison table is declared early so that it floats to the top of page 2.
\nocite{wang2019multitrain,liu2020robust,yang2020biobjective,huang2021integrated,huang2023twostage,liao2020realtime,he2020charging,hu2022joint,li2023cooperative,duan2023integrated,bi2024realtime,wu2022multidepot,li2023trajectory,schulman2017ppo,holland1975adaptation,kennedy1995pso,storn1997de}

\begin{table*}[t]
\centering
\caption{Comparison with representative related work. A dash means that the cited study does not focus on that feature. DRL, SC, LTO, PPO, and MPC denote deep reinforcement learning, supercapacitor, lithium-titanate battery, proximal policy optimization, and model predictive control, respectively.}
\label{tab:positioning}
\scriptsize
\setlength{\tabcolsep}{2.5pt}
\begin{tabular}{p{2.5cm}cccccccc}
\toprule
Work & Scope & Time & Decisions & Joint svc.+chg. & Het. storage & Peak/budget & Feedback & Real/calib. data \\
\midrule
Ghaviha \textit{et al.} \cite{ghaviha2019speed} & Train & trajectory & speed & -- & battery & -- & -- & calibrated \\
Wang and Goverde \cite{wang2019multitrain} & multi-train & timetable & trajectories & -- & -- & -- & -- & operational \\
Huang \textit{et al.} \cite{huang2021integrated} & multi-line & timetable & timetable & -- & -- & -- & -- & case data \\
Hu \textit{et al.} \cite{hu2022joint} & bus network & schedule & location+charging & yes & battery & grid/uncert. & -- & case data \\
Duan \textit{et al.} \cite{duan2023integrated} & bus network & schedule & service+charging & yes & battery & charging plan & -- & case data \\
Bi \textit{et al.} \cite{bi2024realtime} & bus line & real time & flash charging & joint & battery & charger & DRL & case data \\
This work & two tram lines & events & dwell+power & yes & SC+LTO & shared soft budget & PPO/MPC & public+assumed \\
\bottomrule
\end{tabular}
\end{table*}

Multi-train trajectory and timetable studies coordinate running times, headways, and regenerative-energy exchange \cite{wang2019multitrain,liu2020robust,yang2020biobjective}. Other studies extend timetable optimization to connected lines and two-stage urban-rail planning \cite{huang2021integrated,huang2023twostage}, or use deep learning to reschedule trains in real time under random disturbances \cite{liao2020realtime}. This research area mainly controls train movement or the use of regenerative energy; it does not combine arrival-based stop and charging decisions, different onboard storage technologies, and a soft budget shared across lines.

Closely related electric-bus studies jointly optimize charging times, charger locations, vehicle schedules, and timetable changes for fast or opportunity charging \cite{he2020charging,hu2022joint,li2023cooperative,duan2023integrated}. Deep reinforcement learning has also been used for real-time flash-charging control \cite{bi2024realtime}, and multi-depot scheduling models have included power-grid conditions \cite{wu2022multidepot}. These studies demonstrate the value of coordinating service and energy decisions, but depot charging, charger placement, and bus scheduling differ from charging a catenary-free tram at each station.

Taken together, these studies leave unresolved how multiple catenary-free tram lines with heterogeneous onboard storage should coordinate asynchronous stop and charging decisions under a shared demand budget. Treating each line independently cannot fully account for cross-line charging overlap or make use of capacity that is temporarily idle on another service. Table~\ref{tab:positioning} positions this problem against representative related work.

To address this gap, this paper proposes an event-driven deep reinforcement learning method for multi-line catenary-free tramways. A discrete-event simulation captures asynchronous vehicle arrivals and overlapping charging across lines. At each arrival, the online decision-making agent uses the arriving vehicle's state and ongoing charging activity on both lines to jointly determine its stop duration and charging power under the shared demand budget. The case considers two catenary-free tram lines with different onboard storage systems operating during the same simulated morning peak. It is based on Guangzhou Haizhu Tram Line 1 (\thz), which charges onboard supercapacitors at stations, and Guangzhou Huangpu Tram Line 1 (\thp), which uses both supercapacitors and lithium-titanate batteries \cite{haizhu_wiki,haizhu_baike,huangpu_wiki,huangpu_sasac}. The two services are managed under one operator-level demand plan, so their simultaneous grid-side charging powers are added and compared with one shared budget. This penalized supply-planning or purchased-demand target may be exceeded temporarily, while its available capacity remains dynamically shared between the lines.

This paper makes three contributions. First, to the best of our knowledge, it is the first to formulate integrated stop-duration and charging-power control for multiple catenary-free tram lines. The formulation couples passenger-service decisions with charging requirements across asynchronous tram operations and heterogeneous onboard energy-storage systems under an operator-level aggregate power-demand constraint, extending stop--charging coordination from a single line to a multi-line setting. Second, the paper develops an event-driven decision-making framework based on a unified two-line discrete-event simulation and Markov decision process (MDP) formulation. Arrival events from both lines are processed in one event sequence, allowing each decision to account for vehicle and passenger states, traction-energy requirements, operating disturbances, and ongoing charging activities across the two lines. Third, the paper develops an online decision-making agent based on deep reinforcement learning (DRL) for the proposed multi-line environment. Through problem-specific designs of the observation space, continuous action space, and reward function, the agent adaptively determines stop duration and charging power at each arrival event and coordinates transport service and energy demand across the two lines.

\section{Multi-Line Catenary-Free Tramway System Modeling}
\subsection{Dual-Line System and Heterogeneous Storage Configuration}
Evaluating the online agent's arrival-by-arrival decisions requires more than static route and vehicle data: each dwell decision changes the departure time, passenger queue, onboard energy, and charging overlap encountered by later vehicles. We therefore construct a unified discrete-event simulator. Let $\mathcal{L}$ denote the set of lines and, for each line $\ell\in\mathcal{L}$, let $\mathcal{V}_{\ell}$ and $\mathcal{N}_{\ell}$ denote its vehicle and platform sets. A line may use supercapacitor-only storage or a hybrid supercapacitor--battery system. Route, fleet, and storage parameters are simulator inputs; their case-study values are reported in Section~IV.

Vehicles on all lines circulate over a finite simulation horizon $[t_0,t_f]$. Initial vehicle positions, fleet sizes, section running times, selected dwell times, turnaround times, and operating disturbances jointly generate the realized arrival sequence and headways. Control actions consequently affect later headways, passenger accumulation, and cross-line charging overlap through their effects on subsequent arrival times. One complete simulation over $[t_0,t_f]$ is called an \emph{episode}.

\subsection{Event-Driven Operations}
Control is required when a vehicle arrives, while the resulting dwell and charging intervals continue between arrivals. The simulator therefore uses a discrete-event mechanism triggered by vehicle arrivals. Each event $k$ corresponds to tram $i$ arriving at station platform $n$ at time $t_k$. Arrival events from both lines enter one global priority queue and are processed in event-time order; simultaneous events are ordered by line code, station identifier, and tram identifier. After each station service is completed, the simulator updates the vehicle, platform, storage, and operating states and then enqueues that vehicle's next arrival. The two lines therefore evolve on one simulation clock despite their different fleets and operating cycles.

A station stop creates a charging interval $[t_k,t_k+d_k)$, where the dwell time $d_k$ is determined by the joint decision model in Section~III. This half-open interval includes its start but excludes its end, so a vehicle that finishes charging at time $t$ no longer contributes to the power observed at $t$. Conversely, when simultaneous arrivals are processed sequentially, each later event recognizes the charging load established by the preceding events. This mechanism represents cross-line charging overlap and aggregate power in continuous physical time.

\subsection{Running-Time and Passenger Dynamics}
Both the selected dwell time and stochastic running delays shift subsequent arrivals and can therefore change which charging intervals overlap. The simulator updates the departure and next-arrival times as
\begin{align}
    t^{\mathrm{dep}}_k &= t_k + d_k, \\
    t^{\mathrm{next}}_k &= t^{\mathrm{dep}}_k + T^{\mathrm{run}}_{k},
\end{align}
where event $k$ represents vehicle $i$ arriving at the current platform at time $t_k$; $d_k$ is its dwell time; $t^{\mathrm{dep}}_k$ is its departure time; $t^{\mathrm{next}}_k$ is the same vehicle's next-platform arrival time generated at event $k$ and inserted into the global queue, and it may differ from the time of global event $k+1$; and $T^{\mathrm{run}}_k$ is its realized running time over the intervening section. All these time variables are measured in seconds.

Let $\overline{T}^{\mathrm{run}}_k$ be the planned section running time and $z_k$ a delay indicator. The stochastic running-time transition is
\begin{equation}
 T^{\mathrm{run}}_k=\overline{T}^{\mathrm{run}}_k
 \left[1+(\mu_{\mathrm d}-1)z_k\right],
 \qquad z_k\sim\operatorname{Bernoulli}(p_{\mathrm d}).
\end{equation}
Here, $p_{\mathrm d}\in[0,1]$ is the probability that a section is delayed and $\mu_{\mathrm d}\geq1$ is the multiplier applied to its planned running time when a delay occurs. Thus, $z_k=0$ gives the planned time, whereas $z_k=1$ gives $\mu_{\mathrm d}\overline{T}^{\mathrm{run}}_k$. This equation defines the general transition mechanism; the parameter combinations and random-number protocol used in the robustness experiments are specified in Section~IV.

Dwell time is also the time available for passenger service, so shortening a stop to reduce charging overlap can leave more people waiting. Let $\mathcal{N}=\bigcup_{\ell\in\mathcal{L}}\mathcal{N}_{\ell}$ be the set of all platforms and let $Q_n(t)$ be the number of passengers waiting at platform $n\in\mathcal{N}$ at time $t$. Total passenger waiting time during an episode is
\begin{equation}
 W=\sum_{n\in\mathcal{N}}\int_{t_0}^{t_f}Q_n(t)\,\frac{dt}{60},
 \qquad [W]=\text{passenger-minutes}.
\end{equation}
Here, $t_0$ and $t_f$ are the episode start and end times. The integral gives passenger-seconds, which are divided by 60 to obtain passenger-minutes. Between events, $Q_n(t)$ increases according to the specified arrival-rate curve. When a tram arrives, passengers alight first. Passengers who are already waiting can then board, subject to the remaining vehicle capacity and the boarding limit $d_k/\delta$, where $\delta$ is the assumed boarding-service time in seconds per passenger; hence, $d_k/\delta$ is the maximum number of boardings allowed by the dwell time. Passengers who reach the platform during the stop wait for the next tram. This simplifying rule defines passenger service within the simulation. The simulator calculates the area under the piecewise-linear queue curve exactly for every interval between events and during each stop. It applies the same rule when updating the state and calculating waiting time, passengers left at the platform, and passengers unable to board.

\subsection{Charging and Heterogeneous Storage Dynamics}
Charging decisions can be judged feasible only against the energy required after departure, so each action is evaluated with a section-level energy-demand model. The model constructs a triangular or trapezoidal speed profile for each section \cite{ghaviha2019speed} and integrates the acceleration, cruising, and braking phases in SI units. It accounts for traction losses during acceleration and cruising and credits energy recovered through regenerative braking. Auxiliary energy depends on running time and the number of passengers onboard:
\begin{equation}
    W^{\mathrm{aux}}_k = \left(P_{\mathrm{ic}} + P_{\mathrm{oc}} q_{i,k}\right)\frac{T^{\mathrm{run}}_k}{3600}.
\end{equation}
Here, $W^{\mathrm{aux}}_k$ is the auxiliary energy over the section following event $k$ in kWh; $P_{\mathrm{ic}}$ is passenger-independent auxiliary power in kW; $P_{\mathrm{oc}}$ is additional auxiliary power in kW per onboard passenger; $q_{i,k}$ is the passenger count onboard vehicle $i$ over that section; and $T^{\mathrm{run}}_k/3600$ converts seconds to hours. The total section demand is $W_k = W^{\mathrm{trac}}_k + W^{\mathrm{aux}}_k$, where $W_k$, $W^{\mathrm{trac}}_k$, and $W^{\mathrm{aux}}_k$ are the total, traction, and auxiliary energy in kWh.

At event $k$, let $P_k$ denote the grid-side charging power assigned to the arriving vehicle after the feasibility reductions formalized in Section~III. The model assumes that this power remains constant throughout the charging interval $[t_k,t_k+d_k)$ and neglects connection and disconnection delays. Hence, the grid-side energy delivered during the stop is
\begin{equation}
 E^{\mathrm{grid}}_k
 =\frac{1}{3600}\int_{t_k}^{t_k+d_k}P_k\,dt
 =\frac{P_kd_k}{3600},
\end{equation}
where $P_k$ is in kW, $d_k$ is in seconds, and division by 3600 converts the result to kWh. This grid-side energy is converted into stored energy using the efficiency of each storage component. For a vehicle with only a supercapacitor, charging is limited by its unused capacity. For a hybrid vehicle, a fixed onboard energy-management system (EMS) rule charges the supercapacitor first, within its state-of-charge limits. Any remaining feasible energy goes to the battery, subject to its state-of-charge and C-rate limits; the C-rate is the charging or discharging power normalized by the storage component's rated energy capacity. During traction, the vehicle also draws from the supercapacitor first and uses the battery for any remaining demand. Energy that neither component can supply is recorded as unmet traction energy. This fixed priority rule determines the power division between the two components.

\section{Proposed Event-Driven Dwell--Recharge Optimization Method}
\subsection{Event-Driven Decision Problem}
Based on the system dynamics in Section~II, the two-line coordination problem is formulated as a finite-horizon event-driven Markov decision process (MDP). The simulator's full state $s_k$ contains the global event queue, all vehicle and platform states, active charging intervals, travel directions, and cumulative costs. A new decision epoch occurs whenever a vehicle on either line arrives. After the agent selects the dwell and charging decisions for that vehicle, the environment completes passenger service, updates onboard storage, simulates the next section movement, and advances to the next arrival in the global queue. Each decision interval is the elapsed physical time between successive arrival events.

At the level of the full state, $s_k$ and the selected action determine the next-state distribution. The policy, however, receives only the observation $o_k$ defined below. The online decision-making agent consequently faces a partially observable problem. Its objective is to coordinate asynchronous decisions across both lines under the shared power budget, reducing overload and unnecessary charging while limiting unmet traction energy and passenger waiting.

\subsection{Observation and Action Spaces}
The policy cannot use every internal simulator variable in online operation, so it acts on a compact description of the arriving vehicle, its platform, and current charging activity. Before selecting an action, the agent receives these 13 observed variables:
\begin{align}
 o_k = (&t_k,i,n,h_k,E^{\mathrm{sc}}_{i,k},E^{\mathrm{bat}}_{i,k},
 q_{i,k},w_{n,k},\nonumber\\
 &P^{\mathrm{tot}}_k,P^{\mathrm{own}}_k,N^{\mathrm{ch}}_k,
 E^{\mathrm{rem}}_k,\ell_k),
\end{align}
Here, $o_k$ is the policy observation at event $k$; $i$ and $n$ identify the arriving vehicle and platform, and $\ell_k$ identifies the vehicle's line. The arrival time $t_k$ and the time $h_k$ since the previous departure on the same line, at the same platform, and in the same direction are measured in seconds. $E^{\mathrm{sc}}_{i,k}$ and $E^{\mathrm{bat}}_{i,k}$ are the energies stored in vehicle $i$'s supercapacitor and battery at event $k$, in kWh; $E^{\mathrm{bat}}_{i,k}=0$ for a vehicle without a battery. The variables $q_{i,k}$ and $w_{n,k}$ are the onboard passenger count and the number waiting at platform $n$. Before the current action is assigned, $P^{\mathrm{tot}}_k$ is the active charging power across both lines and $P^{\mathrm{own}}_k$ is the portion on the arriving vehicle's line, both in kW. $N^{\mathrm{ch}}_k$ is the number of vehicles currently charging, and $E^{\mathrm{rem}}_k$ is the grid energy still scheduled for those vehicles in kWh. Vehicle and platform identifiers are defined separately within each line.

At event $k$, the agent outputs the two-dimensional continuous action
\begin{equation}
    a_k = (d_k, \widehat{P}_k),
\end{equation}
where $d_k$ is the dwell time and $\widehat{P}_k$ is the requested grid-side charging power. The general action bounds are
\begin{equation}
    d^{\min} \le d_k \le d^{\max}, \qquad
    0 \le \widehat{P}_k \le P^{\max}.
\end{equation}
Here, $d^{\min}$ and $d^{\max}$ are the minimum and maximum allowed dwell times, and $P^{\max}$ is the largest power the agent may request. The simulator maps $\widehat{P}_k$ to an applied power $P_k\leq\widehat{P}_k$ by enforcing station power, remaining storage capacity, charging efficiency, state-of-charge, and C-rate limits. A platform without charging equipment forces $P_k=0$; a request below the no-charge threshold $P^{\mathrm{off}}$ also gives $P_k=0$. Section~IV provides the case-study values of these bounds and thresholds.

For an active interval $r$ with actual charging power $P_r$, start time $\tau_r$, and stop duration $d_r$, the remaining grid-side energy at event time $t$ is
\begin{equation}
 E^{\mathrm{rem}}(t)=\sum_{r\in\mathcal{A}(t)}
 \frac{P_r}{3600}\left[d_r-\min\{\max(t-\tau_r,0),d_r\}\right],
\end{equation}
where $r$ indexes charging intervals, $\mathcal{A}(t)$ is the set of intervals active at time $t$, and the bracketed term is the number of charging seconds remaining in interval $r$. The power $P_r$ is in kW, while $\tau_r$, $d_r$, and $t$ are in seconds; division by 3600 therefore gives $E^{\mathrm{rem}}(t)$ in kWh. The quantity $E^{\mathrm{rem}}(t)$ is the grid energy still required by the current charging schedules, and the observed variable $E^{\mathrm{rem}}_k$ is $E^{\mathrm{rem}}(t_k)$.

\subsection{Shared Power Budget and Cost Functions}
The budget $B$ is a penalized planning target representing either a contractual demand limit or the power capacity purchased by the operator. Temporary exceedance of $B$ incurs a penalty. Because overlapping charging can exceed the budget by different amounts and for different durations, a single sampled power value or peak does not fully describe the violation. Let $\mathcal{I}$ contain the consecutive intervals formed by all charging start and end times in $[t_0,t_f)$. Total charging power is constant within each interval $I\in\mathcal{I}$. The shared-budget overload measure and cost therefore integrate the squared excess over physical time:
\begin{align}
 O_{\mathrm{shared}}(B)&=\sum_{I\in\mathcal{I}}\frac{|I|}{3600}
 [P^{\mathrm{tot}}_I-B]_+^2,\\
 C^{\mathrm{ol}}_{\mathrm{shared}}&=c_oO_{\mathrm{shared}}(B),
\end{align}
where $O_{\mathrm{shared}}(B)$ is the shared-budget overload measure, $P^{\mathrm{tot}}_I$ is the aggregate charging power during interval $I$, $[x]_+=\max(0,x)$, and $|I|$ is the interval duration in seconds. Thus, $O_{\mathrm{shared}}$ is measured in $\mathrm{kW}^2\mathrm{h}$. The coefficient $c_o$ is the overload-cost weight, and $C^{\mathrm{ol}}_{\mathrm{shared}}$ is the resulting cost. Squaring the excess gives larger exceedances greater weight while the interval length retains their duration. To compare flexible sharing with a fixed pre-allocation, the split-budget case uses
\begin{align}
 O_{\mathrm{split}}(B)&=\sum_{\ell\in\mathcal{L}}\sum_{I\in\mathcal{I}}
 \frac{|I|}{3600}[P_{\ell,I}-B_\ell]_+^2,\\
 C^{\mathrm{ol}}_{\mathrm{split}}&=c_oO_{\mathrm{split}}(B),
\end{align}
where $O_{\mathrm{split}}(B)$ is the split-budget overload measure, $\ell$ indexes lines, $P_{\ell,I}$ is the charging power of line $\ell$ during interval $I$, and $B_\ell$ is its allocated budget, with $\sum_{\ell\in\mathcal{L}}B_\ell=B$. The quantity $C^{\mathrm{ol}}_{\mathrm{split}}$ is the corresponding cost. Section~IV specifies the budget allocation used in the experiments. Thus, overload is the physical-time integral of squared excess power.

Overload alone would favor avoiding charging, even when energy is needed for the next sections. The objective therefore also includes regular charging energy, unmet traction energy, and passenger waiting. A separate \emph{anti-cycling excess} term represents cumulative charging substantially above cumulative use. The regular charging cost is
\begin{equation}
 C^{\mathrm{ch}}_k = c_e E^{\mathrm{ch}}_k,
\end{equation}
where $E^{\mathrm{ch}}_k$ is the energy added to onboard storage at event $k$ in kWh, $c_e$ is the regular charging-cost weight, and $C^{\mathrm{ch}}_k$ is the event charging cost. Let $\mathcal{J}_i$ be the set of storage components on vehicle $i$. The quantities $Q^{\mathrm{ch}}_{i,j,k}$ and $Q^{\mathrm{cons}}_{i,j,k}$ are the cumulative energy charged into and consumed from component $j\in\mathcal{J}_i$ up to event $k$, both in kWh. Their cumulative excess is
\begin{equation}
 X_{i,j,k}=\max\!\left(0,Q^{\mathrm{ch}}_{i,j,k}
 -1.25Q^{\mathrm{cons}}_{i,j,k}\right).
\end{equation}
At event $k$, the anti-cycling cost uses the signed change in this cumulative excess:
\begin{equation}
 C^{\mathrm{oc}}_k=c_p\sum_{j\in\mathcal{J}_i}
 \left(X_{i,j,k}-X_{i,j,k^-}\right),
\end{equation}
where $X_{i,j,k}$ is measured in kWh, $k^-$ denotes the instant immediately before event $k$, $c_p$ is the anti-cycling cost weight, and $C^{\mathrm{oc}}_k$ is the event anti-cycling cost. The factor 1.25 is a fixed design margin that discourages unnecessary charge--discharge cycles. When later energy use reduces the cumulative excess, a negative event cost $C^{\mathrm{oc}}_k$ cancels an equal portion of the earlier penalty. Over a complete episode, the increments sum to the nonnegative quantity $c_p\sum_i\sum_{j\in\mathcal{J}_i}X_{i,j,K_i}\geq0$, where $K_i$ is the final event index for vehicle $i$. After accounting for efficiency and enforcing the state-of-charge and C-rate limits of both storage components, the hybrid EMS reports the traction energy still unmet at event $k$ as $D_k\geq0$ in kWh. At episode level,
\begin{equation}
 C^{\mathrm{em}}=c_m\sum_kD_k,\qquad C^{\mathrm{wt}}=c_w W.
\end{equation}
Here, $c_m$ and $c_w$ are the unmet-traction-energy and passenger-waiting cost weights, while $C^{\mathrm{em}}$ and $C^{\mathrm{wt}}$ are the corresponding episode costs. The event-level unmet-energy cost is $C^{\mathrm{em}}_k=c_mD_k$.

\subsection{Optimization Objective and Reward Function}
Charging, anti-cycling, and unmet-energy costs are updated when their corresponding events occur. Overload and passenger-waiting costs are integrated over physical time and assigned incrementally to adjacent decision events. The total cost and reward at event $k$ are
\begin{equation}
 C_k = C^{\mathrm{ch}}_k + C^{\mathrm{oc}}_k + \Delta C^{\mathrm{ol}}_k
 + C^{\mathrm{em}}_k + \Delta C^{\mathrm{wt}}_k,
 \qquad R_k=-C_k,
\end{equation}
where $C_k$ and $R_k$ are the total cost and reward at event $k$, and $\Delta C^{\mathrm{ol}}_k$ and $\Delta C^{\mathrm{wt}}_k$ are the overload and passenger-waiting costs accumulated since the preceding decision epoch. These incremental terms divide the episode-level integrals among events without counting any interval twice. The learning policy $\pi$, implemented using proximal policy optimization (PPO), maximizes the expected event-index-discounted return:
\begin{equation}
 \max_{\pi}\ J_{\gamma}(\pi)
 =\mathbb{E}_{\pi}\!\left[\sum_{k=0}^{K-1}\gamma^k R_k\right].
\end{equation}
Here, $J_{\gamma}(\pi)$ is the expected discounted return of policy $\pi$, $\mathbb{E}_{\pi}$ takes the expectation over policy actions and stochastic running-time disturbances, $\gamma$ is the event-level discount factor, $K$ is the number of decision events in the episode, and the exponent $k$ is the decision-event index.

The coefficients $c_e,c_o,c_m,c_p,$ and $c_w$ are fixed, unit-carrying design weights expressed in arbitrary objective units. Their implemented values and the protocol used to vary them are reported in Section~IV.

\subsection{PPO-Based Solution}
Each arrival requires an online decision from a partial observation, and both control variables are continuous. The solution therefore uses a stochastic feedback policy and adopts proximal policy optimization (PPO), whose clipped importance-ratio objective permits repeated updates from collected trajectories while reducing the optimization incentive for large probability-ratio changes \cite{schulman2017ppo}. The actor $\pi_\theta$, with parameters $\theta$, defines a factorized Gaussian distribution over an unconstrained latent action $u_k\in\mathbb{R}^2$:
\begin{equation}
 u_k\sim\mathcal{N}\!\left(\mu_\theta(o_k),
 \operatorname{diag}(\sigma_\theta^2)\right).
\end{equation}
Here, $\mu_\theta(o_k)\in\mathbb{R}^2$ is the observation-dependent mean, $\sigma_\theta\in\mathbb{R}_{>0}^2$ is the positive standard-deviation vector parameterized by the actor, and the diagonal covariance makes the two latent components conditionally independent. The critic $V_\phi$, with parameters $\phi$, estimates the discounted return from $o_k$.

Let $a^{\min}=(d^{\min},0)^\top$, $a^{\max}=(d^{\max},P^{\max})^\top$, and $\mathbf{1}=(1,1)^\top$. The latent sample is squashed and mapped componentwise to the admissible box ranges for dwell time and requested power:
\begin{equation}
 a_k=a^{\min}+\frac{\tanh(u_k)+\mathbf{1}}{2}
 \odot(a^{\max}-a^{\min}),
\end{equation}
where $\odot$ denotes elementwise multiplication. This construction produces the bounded policy proposal $a_k=(d_k,\widehat{P}_k)$. The simulator then maps the requested power to the physically applied power using the storage and charging-equipment limits described in Section~III-B.

During training, a \emph{rollout} is the complete event sequence collected in one episode. The actor and critic parameters are held fixed at $\theta_{\mathrm{old}}$ and $\phi_{\mathrm{old}}$ while a rollout is generated; $\pi_{\theta_{\mathrm{old}}}$ is therefore the behavior policy that generated its actions and stored log probabilities.

The probability density in the bounded action coordinates follows from the change-of-variables rule. For action dimension $j\in\{1,2\}$,
\begin{align}
 \log\pi_\theta(a_k\mid o_k)
 =&\ \log\mathcal{N}\!\left(u_k;\mu_\theta(o_k),
 \operatorname{diag}(\sigma_\theta^2)\right) \nonumber\\
 &-\sum_{j=1}^{2}\log\!\left(1-\tanh^2(u_{k,j})\right) \nonumber\\
 &-\sum_{j=1}^{2}\log\!\left(\frac{a_j^{\max}-a_j^{\min}}{2}\right),
\end{align}
where $u_{k,j}$, $a_j^{\min}$, and $a_j^{\max}$ are the $j$th components of the corresponding vectors. In numerical evaluation, each factor $1-\tanh^2(u_{k,j})$ is lower-bounded by $\varepsilon_{\mathrm{num}}>0$ before taking its logarithm. The current and behavior policies use the same fixed transformation, so these Jacobian terms cancel in their importance ratio; the expression above defines the transformed density used for the bounded-action log probability and entropy.

After a rollout is collected, generalized advantage estimation (GAE) combines temporal-difference residuals over multiple future events to estimate the policy advantage with a controllable bias--variance trade-off \cite{schulman2015gae}. With $V_{\phi_{\mathrm{old}}}(o_K)=0$ at the terminal event,
\begin{align}
 \delta_k={}&R_k+\gamma V_{\phi_{\mathrm{old}}}(o_{k+1})
 -V_{\phi_{\mathrm{old}}}(o_k),\\
 \widehat{A}_k={}&\sum_{l=0}^{K-k-1}
 (\gamma\lambda_{\mathrm{GAE}})^l\delta_{k+l},\qquad
 \widehat{V}_k=\widehat{A}_k+V_{\phi_{\mathrm{old}}}(o_k).
\end{align}
Here, $\delta_k$ is the temporal-difference residual, $l$ is the nonnegative future-event offset, $\lambda_{\mathrm{GAE}}$ controls the bias--variance trade-off, $\widehat{A}_k$ is the advantage estimate, and $\widehat{V}_k$ is the critic target.

Let the importance ratio computed from the transformed-action probabilities be
\begin{equation}
 \rho_k(\theta)=
 \frac{\pi_\theta(a_k\mid o_k)}
 {\pi_{\theta_{\mathrm{old}}}(a_k\mid o_k)}.
\end{equation}
The actor maximizes the clipped surrogate objective with an entropy bonus:
\begin{align}
 J_{\mathrm{actor}}(\theta)=\mathbb{E}_k\big[&
 \min\{\rho_k(\theta)\widehat{A}_k,
 \operatorname{clip}(\rho_k(\theta),1-\epsilon,1+\epsilon)
 \widehat{A}_k\}\nonumber\\
 &+\beta_H\mathcal{H}(\pi_\theta(\cdot\mid o_k))\big],
\end{align}
where $\mathbb{E}_k$ is the empirical average over events stored in the rollout; $\operatorname{clip}$ truncates its first argument to the interval specified by its last two arguments; $\epsilon$ is the clipping threshold; and $\beta_H$ is the entropy coefficient. The differential entropy of the transformed bounded-action policy is $\mathcal{H}(\pi_\theta(\cdot\mid o_k))=-\mathbb{E}_{a\sim\pi_\theta(\cdot\mid o_k)}[\log\pi_\theta(a\mid o_k)]$. Its coefficient decays linearly during training, assigning more weight to stochastic exploration early in training and progressively reducing this exploration incentive. The critic is trained by minimizing
\begin{equation}
 L_V(\phi)=\frac{1}{2}\mathbb{E}_k
 \left[\left(V_\phi(o_k)-\widehat{V}_k\right)^2\right].
\end{equation}

Each rollout stores observations, latent and bounded actions, rewards, behavior-policy log probabilities, and value estimates. PPO partitions these event samples into shuffled minibatches and reuses them for multiple update epochs, performing gradient ascent on $J_{\mathrm{actor}}$ and gradient descent on $L_V$ before collecting a new rollout. During deterministic evaluation, $u_k=\mu_\theta(o_k)$ is passed through the same hyperbolic-tangent and affine transformations. The network architecture, optimizer, training hyperparameters, entropy schedule, and evaluation settings are reported in Section~IV-D.

\section{Case Study and Experimental Setup}
\subsection{Data Sources and Assumptions}
The case-study data fall into three groups: public facts, derived values, and explicit assumptions. Public facts include line identity, route length, number of stations, fleet size, vehicle capacity, and storage technology when available. When exact station locations are unavailable, station spacing is assumed to be uniform. Section running times are derived from route length and published average speed, and supercapacitor energy is derived from assumed capacitance and voltage. Other assumptions include station-level passenger arrival profiles, boarding-service time, some storage limits, auxiliary power, and the shared power-demand budgets. These public facts, derived values, and explicit assumptions define the model inputs and the scope of the study.

The \thz{} route has 11 stations over 7.7 km, seven vehicles, and a derived usable supercapacitor capacity of 4.734 kWh. The \thp{} route has 19 stations over 14.3 km, 11 vehicles during peak periods, a supercapacitor capacity of approximately 4.66 kWh, and an assumed usable lithium-titanate battery capacity of 8.0 kWh. The simulation horizon is 07:30--09:30, so $t_f-t_0=2$ h; at the initial time, vehicles are placed at different positions in their operating cycles. The model assumes that every platform on both lines can charge a vehicle and sets $d^{\min}=10$ s, $d^{\max}=50$ s, $P^{\max}=300$ kW, and $P^{\mathrm{off}}=5$ kW.

\subsection{Experimental Factors}
The experiments vary four dimensions:
\begin{itemize}
    \item decision-making method: w/o-Opt., GA, PSO, DE, MPC, and PPO;
    \item disturbance environment: no disturbance and four robustness environments $E_1$, $E_2$, $E_3$, and $E_4$;
    \item budget value: $B \in \{1600,1400,1200\}$ kW;
    \item accounting rule: shared aggregate budget versus split per-line budget.
\end{itemize}
The no-disturbance environment sets $p_{\mathrm d}=0$. The parameter pairs $(p_{\mathrm d},\mu_{\mathrm d})$ for the robustness environments are $E_1=(0.2,1.10)$, $E_2=(0.5,1.10)$, $E_3=(0.2,1.20)$, and $E_4=(0.5,1.20)$. For each evaluation seed, the two lines use separate pseudorandom streams derived from that seed to generate their delay-indicator sequences. The same seeds are reused across decision-making methods to reduce differences caused by random samples. The main split-budget comparison sets $B_{\thz}=B_{\thp}=B/2$.

Each combination of $B$ and accounting rule is trained or optimized independently. PPO, GA, PSO, DE, and MPC all use the corrected overload integral over physical time and the passenger waiting-time objective. Only results generated with this corrected objective enter the comparison.

\subsection{Compared Methods}
To compare the learned policy with explicit short-term look-ahead, the model predictive control (MPC) baseline repeatedly optimizes a finite horizon of merged arrival events, applies only the first action, and then solves the horizon again at the next event. Its internal predictor uses planned running times, a table of future events on both lines, a simplified energy model, and the assumption that current aggregate power continues in the background. Its short-horizon objective contains charging-energy, overload, and unmet-energy terms together with an approximation of future passenger loads. The full simulator evaluates every action selected by MPC with the same event costs used for all other decision-making methods. Its internal model and future-event information give MPC an explicit information advantage over PPO.

The fixed-operation baseline, labeled ``without optimization'' (w/o-Opt.), applies the same dwell-time and requested-power decisions at every event, after which the simulator enforces the feasibility limits to obtain $P_k$. To compare online feedback with full-episode offline search, the genetic algorithm (GA), particle swarm optimization (PSO), and differential evolution (DE) optimize the vector $x=[d_1,\ldots,d_K,\widehat{P}_1,\ldots,\widehat{P}_K]$ under the same action bounds. Each vector entry has a fixed link to a line code, tram identifier, and visit index $(\ell,i,v)$. The same vector entry, or gene, therefore controls the same vehicle visit in every disturbance run. GA uses crossover and mutation, PSO uses cognitive, social, and inertia terms, and DE uses differential scaling and recombination \cite{holland1975adaptation,kennedy1995pso,storn1997de}. Their numerical settings are reported with the other experimental configurations in this section.

The fixed policy and search-based baselines minimize the undiscounted episode cost $\sum_{k=0}^{K-1}C_k$ under the same cost definition, ensuring a consistent evaluation objective across methods.

\subsection{Algorithm Configuration and Evaluation Protocol}
Table~\ref{tab:algorithm_settings} reports the numerical configurations used to generate the experimental results. During deterministic PPO evaluation, the latent Gaussian mean is passed through the same hyperbolic tangent and affine transformations used in training, and the observation normalizer is frozen at its training statistics.

\begin{table}[t]
\centering
\caption{Algorithm configurations used in the experiments.}
\label{tab:algorithm_settings}
\footnotesize
\setlength{\tabcolsep}{3pt}
\begin{tabular}{lll}
\toprule
Method & Parameter & Value \\
\midrule
PPO & Architecture & $3\times512$, Tanh \\
PPO & $\gamma/\lambda_{\mathrm{GAE}}/\epsilon$ & $0.993/0.98/0.2$ \\
PPO & Optimizer/learning rate & Adam/$10^{-7}$ \\
PPO & Episodes/epochs/minibatch & $2000/64/64$ \\
PPO & Entropy coefficient & $10^{-2}\rightarrow0$ (linear) \\
MPC & Prediction horizon & 8 merged events \\
w/o-Opt. & Dwell/requested power & 30 s/300 kW \\
GA & Crossover/mutation & $0.9/0.1$ \\
PSO & Cognitive/social/inertia & $0.5/0.5/0.8$ \\
DE & Scaling/recombination & $0.5/0.9$ \\
\bottomrule
\end{tabular}
\end{table}

\subsection{Reward-Weight Settings}
All reference comparison runs use the fixed cost weights in Table~\ref{tab:cost_weights}. The same values are used for every decision-making method, disturbance environment, budget, accounting rule, and random seed. These design coefficients are fixed a priori from the reference method configuration and expressed in arbitrary objective units. Their values remain independent of pilot-run scales and evaluation outcomes.

\begin{table}[t]
\centering
\caption{Reference cost weights used in the experiments. ``Cost'' denotes the arbitrary numerical unit of the optimization objective.}
\label{tab:cost_weights}
\footnotesize
\setlength{\tabcolsep}{3pt}
\begin{tabular}{lccc}
\toprule
Cost term & Weight & Value & Unit \\
\midrule
Stored charging energy & $c_e$ & 5 & cost/kWh \\
Power-budget overload & $c_o$ & 10 & cost/(kW$^2$h) \\
Unmet traction energy & $c_m$ & 1000 & cost/kWh \\
Anti-cycling excess & $c_p$ & 750 & cost/kWh \\
Passenger waiting & $c_w$ & 20 & cost/(passenger-min) \\
\bottomrule
\end{tabular}
\end{table}

\subsection{Coefficient Sensitivity Design}
Let $c_r^{\mathrm{ref}}$ denote a reference value in Table~\ref{tab:cost_weights}. The one-factor sensitivity tests set $c_r=\beta_r c_r^{\mathrm{ref}}$ for $r\in\{o,w,m,p\}$ and $\beta_r\in\{0.5,1,2\}$ while holding the other weights at their reference values. A separate trade-off test interprets the dimensionless multiplier ratio $\beta_o/\beta_w$. Every test retrains or re-optimizes each decision-making method with the same seeds and budgets. The results include the passenger-service, unmet-energy, and overload measures defined below.

\subsection{Metrics and Figures}
Each trained or optimized configuration is evaluated with seeds $\{0,1,2,3,4\}$. Passenger-service measures are total waiting time in passenger-minutes, average waiting time per served passenger, passengers still waiting at the end of the simulation, passengers unable to board because of capacity or stop-duration limits, and vehicle load factor, defined as onboard passengers divided by vehicle capacity. Power-budget measures are the exact squared-excess integral in kW$^2$h, overload duration, defined as the total minutes for which the relevant power exceeds its budget, and peak total power in kW. For split accounting, overload duration and peak power are also reported for each line. Unmet traction energy and the individual cost terms are reported separately. A power trace sampled every second may be used for plots, but all overload measures are calculated from the exact charging start and end times.

The comparative study reports the following outputs:
\begin{itemize}
    \item PPO learning curves under $E_1$--$E_4$ and budget configurations;
    \item aggregate power traces showing $P_{\thz}$, $P_{\thp}$, and $P_{\mathrm{tot}}$ with budget reference lines;
    \item trade-off curves comparing shared and split budgets;
    \item cost-component tables for all decision-making methods;
    \item examples of cross-line coordination, such as \thp{} choosing zero charging power during a \thz{} charging peak.
\end{itemize}

\section{Results and Discussion}
This study asks how an operator should divide a common contractual demand limit or purchased power capacity among tram services, and how that choice affects passenger service and the ability of vehicles to meet traction-energy demand. The shared-budget results should therefore be interpreted as an operator-planning comparison.

A further limitation is that the passenger arrival profiles and several storage parameters are explicit modeling assumptions. The experiments therefore represent the documented engineering scenario defined in this paper, with conclusions scoped to the simulated operation of the two Guangzhou tram lines. Evidence for validity comes from consistent model assumptions, reproducible results with fixed random seeds, and sensitivity tests across disturbance and budget settings.

\section{Conclusion}
This paper presents the first event-driven model that jointly controls stop duration and charging power for multiple catenary-free tram lines. The framework combines arrivals at different times, different onboard storage technologies, passenger service, and a shared operator power-demand budget in one discrete-event simulation and reinforcement-learning model. In the Guangzhou two-line case, the online DRL decision-making agent reduces simultaneous charging peaks and budget exceedance while maintaining sufficient onboard energy and passenger-service quality under running-time disturbances. Its overall performance is better than fixed operation, population-based optimization, and receding-horizon control. These results show that adaptive feedback control can help coordinate power demand and real-time operation across several tram lines.

\bibliographystyle{IEEEtran}
\bibliography{references}

\end{document}
